\section{Becas, Intercambios y Agrupaciones}

\subsection{Beneficios y Becas}
A partir del primer año puedes postular a beneficios como PRONABEC (Beca Continuidad o Permanencia) o la Beca Alianza para el Progreso. También existen pasantías para estudiar un periodo académico fuera del país (como intercambios a la Universidad de Chile). El TEFIM (Tercio Estudiantil) suele organizar charlas informativas sobre traslados internos e intercambios.

\subsection{Centros de Estudiantes}
Se recomienda afiliarse a los grupos de investigación de la facultad:
\begin{itemize}
    \item \textbf{CIIM} (Mecánica - M3) | \textbf{CEDIME} (Mecánica Eléctrica - M4) | \textbf{APEIN} (Naval - M5) | \textbf{CEDIM} (Mecatrónica - M6).
    \item \textbf{ACIA:} Asociación Cultural Isaac Asimov (clubes de música, cine, fotografía, lectura, arte o idiomas).
\end{itemize}
